\newpage
\thispagestyle{empty}
\vspace*{0.35\textheight}
\begin{center}
    {\Huge\textbf{CAPÍTULO III}} \\[0.5cm]
    {\Huge\textbf{MARCO APLICATIVO}}
\end{center}

\newpage
\chapter{Marco Aplicativo}

\section{Metodología}

El proyecto adopta CRISP-DM como metodología principal. Esta elección permite organizar de manera sistemática las fases de comprensión del problema, comprensión de los datos, preparación, modelado, evaluación y despliegue.

\section{Arquitectura propuesta}

La arquitectura objetivo se organiza como sigue:

\begin{center}
\texttt{SICOES → Raw → Clean → RAG Dataset → PostgreSQL → pgvector → LangChain → Streamlit}
\end{center}

Esta arquitectura también incorpora una línea base de búsqueda SQL utilizando \texttt{ILIKE}, con el fin de comparar desempeño frente a la búsqueda semántica.

\section{Preparación de datos}

La preparación contempla limpieza de HTML, normalización de columnas, deduplicación, conversión de fechas, construcción de \texttt{document\_id} y generación del campo \texttt{texto\_rag}. El proyecto conserva Parquet como formato principal versionable del corpus, mientras que los CSV se usan sólo como salidas auxiliares de evaluación y trazabilidad.

\section{Modelado y almacenamiento vectorial}

Se emplea un modelo inicial de embeddings multilingüe de la familia Sentence Transformers, variante \texttt{paraphrase-multilingual-MiniLM-L12-v2}, cuyas representaciones de 384 dimensiones se almacenan en PostgreSQL con \texttt{pgvector}. Esta configuración mejora la cobertura para consultas en español sin alterar el esquema actual de almacenamiento vectorial, y permite búsquedas semánticas, filtros estructurados y estrategias híbridas.

\section{Evaluación}

La evaluación se centrará en tres modos de consulta:

\begin{enumerate}
    \item Búsqueda tradicional por palabra clave mediante SQL e \texttt{ILIKE}.
    \item Búsqueda semántica usando similitud vectorial.
    \item Generación de respuestas mediante RAG con contexto recuperado.
\end{enumerate}

Las métricas consideradas incluyen Precision@5, Recall@5, MRR, hit@5 y análisis cualitativo de resultados. Para esta iteración se construyó un conjunto curado de consultas \texttt{dev}, \texttt{val} y \texttt{test}, permitiendo comparar búsquedas keyword, semánticas e híbridas.

Los resultados obtenidos indican que la búsqueda keyword basada en tokens constituye la línea base más fuerte sobre el corpus actual, seguida por una estrategia híbrida con reranking semántico. La búsqueda semántica pura presenta menor desempeño en las consultas curadas, lo que sugiere que el dominio SICOES mantiene una alta dependencia de terminología específica y coincidencia léxica exacta. En consecuencia, la arquitectura final del proyecto adopta retrieval configurable y no presupone la superioridad universal del enfoque semántico puro.
